\section{増減と凹凸}\label{節:増減凹凸}
$f'(c)=0$ は極値の候補を探す条件であり,極値を保証する条件ではない.\,$f(x)=x^3$ は $f'(0)=0$ だが,\,$0$ の前後でも増加し続ける.極大・極小を判定するには,その点の前後で導関数の符号がどう変わるかを調べる.\,$+$ から $-$ に変われば極大,\,$-$ から $+$ に変われば極小である.

導関数をもう一度微分した $f''=(f')'$ を\textbf{第二次導関数}という.\,$f''>0$ なら接線の傾き $f'$ が増加し,グラフは\textbf{下に凸}になる.\,$f''<0$ なら傾きが減少し,グラフは\textbf{上に凸}になる.たとえば $y=x^2$ は下に凸で,\,$y=-x^2$ は上に凸である.

この呼び方を,接線との位置関係から確かめられる.\,$f''>0$ の区間で $a$ を固定し,\,$x>a$ とすると,平均値の定理により
\[
 f(x)-f(a)=f'(c)(x-a)>f'(a)(x-a)
 \qquad(a<c<x)
\]
となる.\,$x<a$ の場合も,\,$x<c<a$ で $f'(c)<f'(a)$ となることと $x-a<0$ を合わせれば,同じ結論を得る.したがって,曲線は接線の上側にある.\,$f''<0$ では不等号が逆になる.

グラフの凹凸が切り替わる点を\textbf{変曲点}という.\,$f''$ が連続な場合,凹凸が切り替わる場所は $f''=0$ となる点から探せるが,\,$f''=0$ だけでは十分でない.\,$y=x^4$ は $x=0$ で $f''=0$ になっても,その前後でともに下に凸である.

\begin{example}{極値と最大・最小}
$f(x)=x^3-3x$ の増減と凹凸を調べ,区間 $[-2,3]$ における最大値・最小値を求めよ.
\end{example}
導関数を因数分解すると,符号を調べやすい.
\[
 f'(x)=3x^2-3=3(x-1)(x+1),\qquad f''(x)=6x.
\]
したがって $x<-1$ で増加,\,$-1<x<1$ で減少,\,$x>1$ で増加する.\,$f(-1)=2$ が極大値,\,$f(1)=-2$ が極小値である.また,\,$x<0$ で上に凸,\,$x>0$ で下に凸だから,原点 $(0,0)$ は変曲点である.
\[
\begin{array}{c|ccccccc}
 x&-2&&-1&&1&&3\\ \hline
 f'(x)&&+&0&-&0&+&\\
 f(x)&-2&\nearrow&2&\searrow&-2&\nearrow&18
\end{array}
\]
閉区間での最大・最小は,極値だけを比較しては決まらない.両端も候補になるため,\,$f(-2)=-2$,\,$f(3)=18$ を加えて比較する.最大値は $18$ で $x=3$ のとき,最小値は $-2$ で $x=-2,1$ のときにとる.

\begin{figure}[htbp]
\centering
\begin{tikzpicture}[x=2.1cm,y=0.8cm,>=Stealth]
 \draw[->] (-2.3,0)--(2.35,0) node[right]{$x$};
 \draw[->] (0,-3.4)--(0,3.5) node[above]{$y$};
 \draw[very thick,domain=-2.1:2.1,samples=120] plot (\x,{\x^3-3*\x});
 \fill (-1,2) circle(1.5pt) node[above left]{$(-1,2)$};
 \fill (1,-2) circle(1.5pt) node[below right]{$(1,-2)$};
 \fill (0,0) circle(1.5pt) node[above right]{変曲点};
 \draw[dashed] (-1,0) node[below]{$-1$}--(-1,2)--(0,2);
 \draw[dashed] (1,0) node[above]{$1$}--(1,-2)--(0,-2);
 \node at (1.1,2.85) {$y=x^3-3x$};
\end{tikzpicture}
\caption{極値は増減の切り替わり,変曲点は凹凸の切り替わりを表す.}
\label{図:三次増減}
\end{figure}

図\ref{図:三次増減}は原点付近を描いたものである.見えている山の高さ $2$ が,指定区間全体の最大値とは限らないことに注意しよう.最大・最小を求めるときは,まず定義域と対象の区間を確かめ,内点の極値候補,微分できない点,含まれる端点を比較する.連続関数でも,区間が閉じていなければ最大・最小が存在しない場合がある.


